% =============================================================================
% Section V — Results
% =============================================================================
\section{Results}
\label{sec:results}

\subsection{Quantitative Performance on Test Set}
Table~\ref{tab:metrics} reports the final evaluation on the 221-window test
set at the calibrated threshold $\tau^{*} = 0.0200$. The AUC-PR of 0.4552
is \textbf{10.5\,\% above the no-skill random baseline} of 0.4118 (equal to
the test set positive rate of 91/221), confirming that the model learns a
meaningful anomaly-scoring function.

\begin{table}[H]
\centering
\caption{Final test-set metrics at $\tau^{*} = 0.0200$.}
\label{tab:metrics}
\begin{tabular}{lc}
\toprule
\textbf{Metric} & \textbf{Value} \\
\midrule
AUC-ROC                   & 0.5011 \\
AUC-PR                    & 0.4552 \\
AUC-PR random baseline    & 0.4118 \\
\midrule
Threshold $\tau^{*}$      & 0.0200 \\
Leak Precision            & 0.4118 \\
Leak Recall               & 1.0000 \\
Leak F1-score             & 0.5833 \\
Accuracy                  & 0.4118 \\
\bottomrule
\end{tabular}
\end{table}

\subsection{Confusion Matrix Analysis}
Figure~\ref{fig:cm_eval} shows the raw-count and row-normalised confusion
matrices at $\tau^{*} = 0.0200$. All 221 test windows are predicted as Leak,
meaning the 91 true leak windows are all captured (recall = 1.0) while all
130 normal windows are false-positives. This recall-first behaviour is a
direct consequence of the very low calibrated threshold: the model assigns
probabilities predominantly above 0.020 to all windows.

\begin{figure}[H]
\centering
\includegraphics[width=\columnwidth]{images/fig_confusion_matrix.jpg}
\caption{Confusion matrix (raw counts and row-normalised) on the test set
at $\tau = 0.0200$. All 221 windows are predicted Leak, yielding perfect
recall and maximum false-alarm rate at this operating point.}
\label{fig:cm_eval}
\end{figure}

\subsection{Ranking Performance (PR/ROC)}
Figure~\ref{fig:prroc_eval} shows the precision-recall and ROC curves on
the test set. The AUC-PR of 0.455 demonstrates meaningful score
separability above the random baseline (AP = 0.41). The steep early rise
in the PR curve indicates that the model correctly assigns the highest
scores to a subset of true leak windows before precision degrades, while
the ROC curve (AUC = 0.501) remains near the diagonal, reflecting the
difficulty of the task on this test split.

\begin{figure}[H]
\centering
\includegraphics[width=\columnwidth]{images/fig_pr_roc_curves.jpg}
\caption{Precision-recall (left) and ROC (right) curves on the test set.
The red dot marks the operating point at $\tau^{*} = 0.020$. The dashed
line indicates the no-skill random baseline.}
\label{fig:prroc_eval}
\end{figure}

\subsection{Threshold Behaviour and Score Distribution}
Figure~\ref{fig:score_dist_eval} shows the predicted probability distributions
for Normal and Leak windows on the test set. Both distributions overlap
substantially in the 0.2--0.8 range, explaining the limited AUC-ROC. However,
Leak windows (n=91) show a marginally higher density at elevated scores,
confirming a directionally correct ranking signal. The calibrated threshold
$\tau^{*} = 0.020$ lies far to the left of both distributions, causing every
window to be predicted Leak.

\begin{figure}[H]
\centering
\includegraphics[width=\columnwidth]{images/fig_score_distribution.jpg}
\caption{Predicted probability distributions for Normal (blue, n=130) and
Leak (red, n=91) windows. The dashed vertical line marks $\tau^{*}=0.020$.
Substantial overlap between classes reflects inherent task difficulty.}
\label{fig:score_dist_eval}
\end{figure}

Figure~\ref{fig:th_sens_eval} illustrates how precision, recall, and F1
vary across the full threshold range. Recall remains at 1.0 for all
thresholds below approximately 0.60, while precision and F1 peak near
$\tau = 0.55$--$0.60$ where the curves intersect.

\begin{figure}[H]
\centering
\includegraphics[width=\columnwidth]{images/fig_threshold_sensitivity.jpg}
\caption{Leak precision, recall, and F1-score as a function of decision
threshold. The dashed vertical line marks the deployed threshold
$\tau^{*} = 0.020$. Peak F1 is achieved near $\tau = 0.58$.}
\label{fig:th_sens_eval}
\end{figure}

